\documentclass[11pt,a4paper]{article}
\usepackage{amsmath,amssymb,amsthm}
\usepackage[margin=2.5cm]{geometry}
\usepackage{hyperref}

\newtheorem{definition}{Definition}[section]
\newtheorem{theorem}[definition]{Theorem}
\newtheorem{proposition}[definition]{Proposition}
\newtheorem{lemma}[definition]{Lemma}
\newtheorem{corollary}[definition]{Corollary}
\newtheorem{conjecture}[definition]{Conjecture}
\newtheorem{remark}[definition]{Remark}

\title{Hyperseed Formalization:\\
  Linguistic Universals as Shadows of Cognitive Structure, Part 2}
\author{ProtomegaTron (automated formalization)}
\date{2026-09-18}

\begin{document}
\maketitle

\begin{abstract}
We formalize the intellectual structure of Ben Goertzel's Substack article
``Linguistic Universals as Shadows of Cognitive Structure, Part 2''
(2026-05-30) within the Hyperseed ontology. The article develops the
categorical formalism that turns the empirical observations of Part 1
into theorem-shaped statements: quantale weakness as a generalized Occam
principle (fiber-minimal explanation), the five linguistic-cognitive
correspondences as transport theorems (fiber morphisms from cognitive
bundle to linguistic bundle), with the factorization correspondence as
the strongest result (an identity theorem = automatic fiber product
preservation). Double-counting in correlation-only explanations is
identified as scalar collapse, the TUG trace quotient as fiber quotient,
and the convergence of independent mathematical lines as evidence for the
fiber-bundle formulation of cognitive architecture.
\end{abstract}

\tableofcontents

%% =================================================================
\section{Quantale weakness as fiber-minimal explanation}
\label{sec:weakness}
%% =================================================================

\begin{theorem}[Generalized Occam --- claims 3--9, 19]
\label{thm:weakness}
The article introduces \emph{quantale weakness} as the selection
principle for the underdetermined inverse problem: given a stable
linguistic pattern, infer the cognitive structure that produced it.

Standard Occam's razor (``shorter description = better'') is
awkward in this context: a short description can be wrong-headed,
and a long description can be the natural one. What you actually
want is the explanation that ``commits to as little as necessary,
that doesn't multiply hidden distinctions, that doesn't represent
the same dependency through several overlapping mechanisms when
one would do.''

Building on Michael Timothy Bennett's set-theoretic weakness
concept, the framework gives this idea an algebraic (quantale) form.
Explanations are compared along five axes simultaneously:
\begin{enumerate}
\item \textbf{Admissible indistinction:} how much structure the
  explanation leaves unresolved (more = better, all else equal).
\item \textbf{Fit:} how well the explanation matches the data.
\item \textbf{Cost:} computational and representational expense.
\item \textbf{Confusion:} cross-context interference introduced.
\item \textbf{Double-counting:} redundant representation of the
  same dependency through multiple pathways.
\end{enumerate}

An explanation wins by dominating competitors in this product
order --- not because someone called it elegant.

The key result: \emph{causal explanations beat correlation-only
explanations as a theorem in the weakness order}. A causal
explanation posits a single genuine mediator (one dependency
represented once). A correlation-only explanation captures the
same regularity through pairwise links, representing the same
dependency many times over. The causal explanation wins on
double-counting, confusion, and usually cost. ``Causal'' is not
a magic word --- it's a structural property that the weakness
order can detect.

In Hyperseed terms, quantale weakness selects the
\emph{fiber-minimal explanation}: the explanation that introduces
the least fiber structure beyond what the data requires.
Unnecessary fiber structure = unnecessary distinctions =
weakness-dominated explanation. The product order on explanations
is the product order on fiber structures: an explanation with
less unnecessary fiber dominates one with more.
\end{theorem}

%% =================================================================
\section{Double-counting as scalar collapse}
\label{sec:double-counting}
%% =================================================================

\begin{proposition}[Redundant pathways --- claim 21]
\label{prop:double-counting}
The article identifies double-counting as the axis on which causal
explanations most clearly dominate correlation-only explanations.
A correlation-only explanation represents the same dependency
through multiple redundant pathways: if A causes B and B causes C,
the correlation-only approach might represent A-B, B-C, and A-C
as three separate dependencies, when in fact the third is entirely
determined by the first two.

In Hyperseed terms, this is \emph{scalar collapse applied to
explanatory structure}: the genuine fiber structure (one mediator,
one chain of dependence) is collapsed into multiple overlapping
scalar summaries (pairwise correlations) that each partially capture
the same underlying structure. The correlations are the scalars;
the causal structure is the fiber. Replacing fiber with scalars
loses the dependency structure and introduces double-counting.

This is the same anti-scalar-collapse principle that appears
throughout the Hyperseed ontology: $(f,c)$ vs.\ scalar truth,
scoped reputation vs.\ scalar trust, essential laterality vs.\
participation entropy. In each case, collapsing structured
information into scalar summaries loses the structure that
determines correct inference.
\end{proposition}

%% =================================================================
\section{The five transport theorems as fiber morphisms}
\label{sec:transport}
%% =================================================================

\begin{theorem}[Externalization functor --- claims 10--15, 20, 23]
\label{thm:transport}
The five linguistic-cognitive correspondences, sharpened in this
article from informal observations to theorem-shaped statements,
are all instances of a single categorical structure: the
\emph{externalization functor} from the cognitive category to the
linguistic category, with quantale weakness as the selection
principle on the inverse.

In Hyperseed terms, each correspondence is a \emph{fiber morphism}
from the cognitive bundle to the linguistic bundle --- a map that
preserves the relevant fiber structure:

\begin{enumerate}
\item \textbf{Factorization (identity theorem).} The strongest
  result. Cognitive updates on disjoint causal modules commute in
  the symmetric monoidal sense. Their linguistic images are
  swap-equivalent in the TUG trace quotient. If cognition factors
  and externalization preserves factorization, the linguistic core
  factors too.

  In Hyperseed terms, this is \emph{automatic fiber product
  preservation}: the externalization functor preserves product
  structure without additional assumptions. This is the strongest
  possible structural result --- an identity, not a conditional.

  The remarkable corollary: ``broad cross-module universals are
  rare'' and ``causal coding bounds catastrophic forgetting'' are
  the \emph{same theorem}. Two apparently independent results ---
  one empirical (linguistics), one theoretical (machine learning)
  --- are revealed as two facets of the same fiber product
  preservation.

\item \textbf{Hierarchical (closure transport).} A closure system
  on the cognitive side pushes forward to a closure system on the
  linguistic side. Fixed points correspond:
  \[
    \text{Fix}(C_{\text{cog}}) \xrightarrow{F}
    \text{Fix}(C_{\text{lang}})
  \]
  Person, number, case, accessibility hierarchies are linguistic
  shadows of cognitive closure structure protected by the
  kernel-shell architecture.

\item \textbf{Word-order (sparse-energy pushforward).} A sparse
  signed-energy controller on the cognitive side produces, after
  marginalizing over non-externalizing variables, a sparse
  signed-energy network at the linguistic level. The pushforward
  preserves sparsity provided the externalization map is local
  enough.

\item \textbf{Graded universalhood (stability-rate transport).}
  The linguistic stability coefficient is the image of an
  underlying cognitive stability coefficient:
  \[
    \sigma_{\text{lang}}(u) = F_*(\sigma_{\text{cog}}(u))
  \]
  Hard-kernel cognitive invariants map to high-stability linguistic
  universals; shell-level cognitive patterns map to lower-stability
  linguistic tendencies.

\item \textbf{Context-indexed (glued-closure transport).} Local
  closure operators on context-specific cognitive substrates glue
  into a global linguistic structure via transition maps. The
  gluing respects the closure structure --- the sheaf condition
  is preserved by the externalization functor.
\end{enumerate}

The unifying insight: these are not five separate observations
stitched together but ``five facets of one categorical structure.''
The externalization functor is a \emph{bundle morphism}: it maps
cognitive fibers to linguistic fibers while preserving the five
structural properties (product, closure, sparsity, stability,
gluing). The five correspondences are the five structural invariants
of this morphism.
\end{theorem}

%% =================================================================
\section{The TUG trace quotient as fiber quotient}
\label{sec:tug}
%% =================================================================

\begin{definition}[Observable equivalence --- claim 22]
\label{def:tug}
The TUG trace quotient identifies derivations that produce the same
observable output. Two derivations that differ in their internal
cognitive steps but produce indistinguishable linguistic output are
identified in the quotient.

In Hyperseed terms, this is a \emph{fiber quotient}: it identifies
sections that project to the same base point. Derivations that
differ in fiber detail (internal cognitive structure) but agree on
observable (base) structure (linguistic output) are identified.
The quotient strips away the fiber information that doesn't
externalize --- the ``hidden'' cognitive structure that leaves no
linguistic trace.

The quotient is what makes the inverse problem (inferring cognition
from language) well-posed: by working in the quotient, you're asking
``what is the simplest cognitive structure consistent with this
linguistic output?'' rather than ``what is the exact cognitive
structure?'' The quantale weakness principle operates in this
quotient space, selecting among equivalence classes of cognitive
explanations.
\end{definition}

%% =================================================================
\section{Convergence as fiber-bundle evidence}
\label{sec:convergence}
%% =================================================================

\begin{proposition}[Independent lines, same structure --- claims 16--18, 25]
\label{prop:convergence}
The article's meta-theoretical claim: the mathematical structure
discovered by linguistic typology (empirical, studying language
across cultures) and the mathematical structure required by
brain-like continual learning (theoretical, designing AI
architectures) converge on the same five features:

\begin{itemize}
\item Closure systems with diachronic Lyapunov flow.
\item Factorization theorems with mediator exceptions.
\item Sparse signed interaction graphs.
\item Graded stability under transition dynamics.
\item Gluing of context-indexed local closure operators.
\end{itemize}

In Hyperseed terms, this convergence is evidence for the
fiber-bundle formulation of cognitive architecture: if two
independent investigations arrive at the same fiber-bundle
structure (decomposition, commutativity, transition functions,
depth, cocycle conditions), that structure is a genuine property
of cognition, not an artifact of either investigation's methodology.

The convergence also validates the Hyperseed ontology itself:
the fiber-bundle vocabulary is not imposed on the data but
emerges naturally from both the empirical and theoretical
investigations. The correspondences are not analogies --- they
are structural identities (in the case of factorization) or
conditional transport results (in the other four cases).

The article's concluding observation --- ``the field of linguistics
has been measuring the architecture of the mind through language
without fully realizing it'' --- is the claim that linguistic
typology is an indirect form of fiber inspection: inspecting the
cognitive fiber through its linguistic projections.
\end{proposition}

%% =================================================================
\section{Cross-references}
%% =================================================================

\begin{remark}[Connections to other formalizations]
\begin{itemize}
\item \textbf{Linguistic Universals Part 1} (2026-05-20): empirical
  foundation and causal-coding architecture (HBCML/ColBaC). This
  article formalizes the correspondences that Part 1 identifies
  informally.
\item \textbf{Linguistic Universals Part 3} (2026-06-08): AGI
  architecture implications and concrete Transformer training
  recipe. Part 3 turns the transport theorems of this article into
  differentiable losses.
\item \textbf{$(f,c)$-lossy} (LTM 5a4d150b): anti-scalar-collapse.
  Double-counting in correlation-only explanations is the same
  scalar collapse that afflicts $(f,c)$-lossy truth fusion.
\item \textbf{Non-invertible 2-cells} (LTM 85c02e32): composition
  structure. The factorization identity theorem is the positive
  version --- when composition works correctly (factorizable
  modules), the fiber product is preserved.
\item \textbf{Avoiding AGI Catastrophe Part 2} (2026-06-10):
  cognitive-cryptographic alignment. The externalization functor
  (cognition $\to$ language) is structurally parallel to the
  cognitive-cryptographic alignment (cognition $\to$ security):
  both are bundle morphisms preserving structural invariants.
\item \textbf{The Orchard Bug} (2026-06-05): cocycle defects.
  The glued-closure transport (C5) is the positive version of
  the Orchard bug: when transition maps respect closure structure,
  the global system is consistent.
\end{itemize}
\end{remark}

\begin{conjecture}[Universality of the externalization functor]
Conjecture: the externalization functor from cognition to language
is an instance of a more general \emph{projection functor} from
any high-dimensional cognitive bundle to any lower-dimensional
observable bundle. The five structural invariants (product, closure,
sparsity, stability, gluing) are preserved by any sufficiently
well-behaved projection, not just the cognition-to-language
projection.

If this conjecture holds, the same five correspondences should
appear in other cognitive-to-observable projections: cognition to
motor behavior, cognition to social interaction, cognition to
artistic expression. Each observable domain would serve as a
different ``shadow'' of the same cognitive architecture, and the
transport theorems would generalize to all of them.
\end{conjecture}

\end{document}
